\section{Using the Model to Answer the Research Question}

\subsection{Policy experiments}

The model is suited to answer the research question through a set of counterfactual simulations. We solve the model separately under three policy regimes:

\begin{enumerate}
    \item \textbf{Baseline} ($\bar{q} = 0$): the household freely assigns leave. Given symmetric preferences and $k_{1,0} = k_{2,0}$, the model predicts that leave is shared equally in the first period --- but once a small wage gap opens (e.g., due to simulation noise or a slight initial asymmetry), the lower earner takes all subsequent leave, generating a self-reinforcing wage gap. This mirrors the empirical pattern in Denmark.

    \item \textbf{Partial quota} ($\bar{q} = 0.1$): fathers must take at least 10\% of total leave. This is close to the Norwegian ``daddy quota'' as originally introduced.

    \item \textbf{Equal sharing} ($\bar{q} = 0.5$): mandatory 50/50 split. A more radical reform, comparable to Iceland's parental leave system.
\end{enumerate}

For each regime, we simulate a panel of $N = 1{,}000$ households and compute:

\begin{itemize}
    \item \textbf{Event-study graph}: average wages of fathers and mothers relative to the year before first birth, as in Figure 1B of \citet{kleven2019children}. Under the baseline, we expect a sharp drop in maternal wages at $t = 0$ and no effect on paternal wages; under equal sharing, both profiles are symmetric.

    \item \textbf{Life-cycle wage gap} $G_t$ (equation \ref{eq:gap}): the gender wage gap as a function of age, for each quota level $\bar{q}$.

    \item \textbf{Labor supply profiles}: average hours $h_{1,t}$ and $h_{2,t}$ over the life cycle, isolating how the quota reshapes each partner's incentive to supply labor.

    \item \textbf{Welfare comparison}: the value function at the beginning of life, $V_0(a_0, k_0, k_0, 0)$, as a measure of lifetime household welfare under each regime. This allows us to ask whether a quota is welfare-improving despite the short-run income cost.
\end{itemize}

\subsection{Identifying the welfare trade-off}

The welfare analysis is particularly informative because the model captures two opposing forces:

\paragraph{Gain: wage convergence.}
By forcing fathers to share the career cost, the quota prevents the divergence in human capital described in equation \eqref{eq:hc}. Over the life cycle, the mother's wage recovers relative to the baseline, increasing her labor supply and household income. This gain is large when $\delta_2$ is high (strong depreciation) and when the time horizon is long.

\paragraph{Cost: short-run income loss.}
If the father is the higher earner at the time of birth (which can be the case even with equal initial capital if random draws of birth timing coincide with differential accumulation), forced leave reduces total household income in that period. This is an income effect that lowers consumption and welfare in the short run.

Whether the gain dominates the cost depends on structural parameters ($\delta_1$, $\delta_2$, $\alpha$, $\phi$) that can be estimated from data. This is a key strength of the structural approach: the model can be used to compute the welfare-maximizing quota level $\bar{q}^\star$.

\section{Data Requirements}

To bring the model to the data and estimate its structural parameters, we would use the following sources:

\paragraph{Wages and earnings over the life cycle.}
Danish administrative register data (IDA) provides matched employer-employee records with annual wages, hours, and labor market status for the full population. The key moments to match are: (i) the wage-age profile by gender, (ii) the child penalty in female earnings (from \citet{kleven2019children}), and (iii) the evolution of the wage gap over the first 10 years after birth.

\paragraph{Leave-taking behavior.}
Statistics Denmark's Barselstatistik records the leave spell duration taken by each parent, including the number of weeks and the benefit rate. These data allow direct calibration of $\bar{L}$ and the baseline distribution of $\lambda_t$ (father's share of leave), which is the key endogenous variable the quota is designed to shift.

\paragraph{Human capital depreciation during leave ($\delta_i$).}
The depreciation parameters $\delta_1$ and $\delta_2$ are the most important structural parameters for the counterfactual exercise. They can be estimated from the event-study regression in \citet{kleven2019children}: the slope of the wage path after birth relative to the pre-birth trend identifies how much human capital is lost per unit of leave taken. Since fathers who voluntarily take leave can be compared to similar fathers who do not, $\delta_1$ can be identified from the (small) wage effect of paternity leave take-up, using policy reforms as instruments \citep{patnaik2019reserving}.

\paragraph{Labor supply and hours.}
The Danish Labor Force Survey (AKU) provides quarterly data on hours worked, employment status, and parental status. These moments are used to calibrate $\beta_0$, $\beta_1$, and $\gamma$ --- the preference parameters governing the disutility of work and its sensitivity to children.

\paragraph{Leave replacement rate.}
Danish law entitles parents to \textit{barselsdagpenge} at a rate of up to 90\% of previous earnings, capped at DKK 4,460 per week (2024). This directly pins down $\phi$ without estimation.

\paragraph{Calibration strategy.}
We propose a simulated method of moments (SMM) approach, as covered in Lecture 2 on structural estimation. The moment conditions include the average wage gap at ages 10, 20, and 30 after labor market entry, the share of leave taken by fathers in the baseline, and the elasticity of female labor supply with respect to the child subsidy (from Assignment 1). Parameters are estimated by minimizing the distance between model-simulated moments and their empirical counterparts.
